\chapter{Introduction}\label{ch:intro}

\section{Introduction}\label{sec:intro}

The models that perform best on tabular data are tree ensembles whose reasoning is not
readable from their parameters~\cite{grinsztajn2022why}, and they are increasingly
deployed where decisions affect people. Regulation has followed: the EU AI Act obliges
providers of high-risk systems to make operation interpretable~\cite{euaiact2024}.
Two post-hoc methods dominate practice. \shap{} attributes a prediction to features
through cooperative-game theory~\cite{lundberg2017unified,lundberg2020local}; \lime{}
fits an interpretable surrogate in the neighbourhood of a single
prediction~\cite{ribeiro2016why}. Their foundations differ, one is deterministic on
trees while the other samples, and their failure modes are
unrelated~\cite{slack2020fooling,kumar2020problems}. A practitioner who runs both and
obtains two different top-five lists has no principled basis for choosing.

Prior comparisons do not settle this because they hold too much fixed: typically one
model family or one dataset, a locally defined fidelity metric, no hypothesis test, no
random-ranking baseline, and no statement of which output space the attributions occupy.
The gap is not a missing algorithm but a missing controlled comparison --- and that is
something an experiment can close.

This paper contributes:

\begin{enumerate}
  \item \textbf{A reproducible multi-axis benchmark.} Two explainers, three rank-fusion
    rules, and a random control across six classifiers and three public clinical
    datasets (\NumCells{} cells, \NumExplained{} instances), with eleven pre-registered
    hypotheses, non-parametric tests, Holm correction within family, effect sizes, and
    bootstrap intervals (Table~\ref{tab:hypotheses}).
  \item \textbf{A three-way decomposition of explanation stability} --- explainer
    randomness, input perturbation, and variance from retraining the model itself
    (Section~\ref{sec:rq2}). Reporting one stability number conflates these, and the
    distinction matters: an explanation can be bit-identical across runs and still
    fragile, because the model beneath it moves when the training sample does.
  \item \textbf{A hybrid consensus method, and an honest account of its limits.} We
    propose \emph{fidelity-weighted rank fusion}: the two explainers are combined in
    proportion to the faithfulness each demonstrates, with the weight fitted on a
    calibration split disjoint from the instances the method is later tested on, so it
    cannot win by construction. It is evaluated against two published fusion baselines ---
    mean-rank (Borda) and reciprocal-rank fusion~\cite{cormack2009reciprocal} --- and
    against a deliberately strict bar: beat the \emph{better} of \shap{} and \lime{} on
    the same cell, not their average, which anything would beat. Explanation aggregation
    itself is prior work~\cite{bhatt2020aggregating} and we claim no priority for the
    idea; what is new is the fidelity weighting, the leakage-safe calibration, and the
    controlled test. The outcome is negative, it was pre-registered as a valid result
    before running, and Section~\ref{sec:rq6} explains the mechanism.
\end{enumerate}

Section~\ref{sec:method} states the protocol, including four methodological traps that
silently invalidate comparisons of this kind. Section~\ref{sec:results} reports one
subsection per research question.


\section{Problem statement}
Two post-hoc explanation methods dominate practice on tabular data, they rest on
different theory, they routinely disagree about the same prediction, and the published
comparisons between them vary too few axes to generalise from. A practitioner choosing
between them has no controlled evidence to consult. This thesis supplies that evidence.

\section{Objectives}
The work is organised around six research questions, each with a measurable outcome:

\begin{enumerate}
  \item[\textbf{O1}] Quantify how far SHAP and LIME agree on feature rankings, and
    whether agreement depends on model family or task difficulty (\RQ{1}).
  \item[\textbf{O2}] Separate explanation stability into explainer randomness, input
    perturbation, and model retraining variance, and measure each (\RQ{2}).
  \item[\textbf{O3}] Measure how faithfully each explanation reflects the model's
    decision function, against a random-ranking control (\RQ{3}).
  \item[\textbf{O4}] Measure the computational cost of each explanation and the
    fidelity that cost buys (\RQ{4}).
  \item[\textbf{O5}] Check top-ranked features against a pre-registered,
    literature-derived reference set (\RQ{5}).
  \item[\textbf{O6}] Test whether rank-level fusion of the two beats the better of them
    (\RQ{6}).
\end{enumerate}

\section{Scope and delimitations}
This is a computer-science methodology study. It covers post-hoc feature-attribution
methods for binary classification on small tabular datasets. It does not cover
counterfactual explainers, rule extraction, attention-based methods, image or text
modalities, deep tabular architectures, or human evaluation of explanation
usefulness. It makes no medical claim of any kind; the three diseases are benchmarks.

\section{Structure of this thesis}
Chapter~\ref{ch:literature} surveys the two methods, their evaluation, their known
failure modes, and prior comparisons. Chapter~\ref{ch:method} states the full protocol.
Chapter~\ref{ch:results} reports results per research question, with all eighteen cells
in Appendix~\ref{app:fullresults}. Chapter~\ref{ch:discussion} interprets them,
Chapter~\ref{ch:limitations} bounds them, and Chapter~\ref{ch:conclusion} closes.
Appendices give the reference feature sets with citations, the recorded environment,
and instructions for reproducing every number.
